\section{Implementation}
\label{sec:impl}

\subsection{The receive path cannot block}
Bytes arrive one at a time in the USART1 receive interrupt and are appended to
a line buffer. The interrupt does no parsing: it tests for the terminator, and
only when a complete line is present does it hand that line to the protocol
layer. Nothing in the receive path waits on anything, so a client that stops
mid-message costs the board nothing but a buffer that the 2000\,ms timeout will
later clear. An overrun is cleared and the byte dropped rather than left to
wedge the receiver, because the parser already recovers at the next newline.

\subsection{Three checks, one write}
\texttt{game\_move()} tests range, then turn, then occupancy, and each check
returns before the board is touched:

\begin{verbatim}
if (index < 0 || index >= GAME_CELLS)
    return MOVE_ERR_RANGE;
if (player != 0 && player != g->turn)
    return MOVE_ERR_NOTURN;
if (g->cells[index] != '.')
    return MOVE_ERR_OCCUPIED;

g->cells[index] = g->turn;
\end{verbatim}

There is exactly one assignment to \texttt{g->cells}, and all three rejection
paths return above it. That structure, rather than an inspection of every path,
is what makes the authority argument in Section~\ref{sec:authority} hold.

\subsection{One core, two targets}
\texttt{server/core/} contains the rules and the frame parser in plain C with
no HAL and no hardware headers. \texttt{server/stm32/port\_stm32.c} is the only
file that knows about the board, and it provides the functions declared in
\texttt{port.h}: bring up the USART, service its interrupt, transmit a line,
drive the LED, read the button and report a millisecond tick. The same core
therefore builds unchanged for the board and for the host simulator in
\texttt{sim/}, so a rule bug is found on a laptop in milliseconds rather than
on the bench.

The board project in \texttt{board/} reaches the shared sources through linked
folders rather than copying them, so there is one copy of every file and the
two builds cannot drift apart. GPIO goes through the HAL because the GPIO
driver ships with the course project, while USART1 is configured from registers
against the reference manual~\cite{rm0091} because the vendored HAL has no UART
driver at all. At 48\,MHz and 115\,200 baud with oversampling by 16 the baud
rate register is the divisor itself, so \texttt{BRR} is 417, which gives
115\,107 baud and an error of 0.08\,per cent.

The release build occupies 4984 bytes of flash and 2720 bytes of RAM, against
the 32\,kB of flash and 8\,kB of RAM that the linker script reserves. The board
used here in fact reports 64\,kB of flash and 8\,kB of RAM when probed, so
neither figure is close to a constraint either way.
